% 小行星卫星（综述）
% license CCBYSA3
% type Wiki

本文根据 CC-BY-SA 协议转载翻译自维基百科\href{https://en.wikipedia.org/wiki/Minor-planet_moon}{相关文章}。

一个小行星卫星是指作为天然卫星环绕一颗小行星运行的天文对象。截至 2022 年 1 月，已有 457 颗已知或被怀疑拥有卫星的小行星。之所以说对小行星卫星（以及一般意义上的双天体系统）的发现十分重要，是因为通过确定它们的轨道，可以估计主星的质量与密度，从而获得通常难以通过其他方式获取的物理性质信息。

在这些卫星中，有若干在相对尺寸上非常巨大：90 号 Antiope、Mors–Somnus 与 Sila–Nunam 的尺寸约为其主星的 95\%；Patroclus–Menoetius、Altjira 与 Lempo–Hiisi 的尺寸约为 90\%，其中 Lempo–Paha 约为 50\%。已知绝对尺寸最大的小行星卫星是冥王星的卫星 Charon，其直径约为冥王星的一半。

此外，在若干遥远的天体周围，人们还发现了多个环系统（参见：Chariklo 与 Chiron 的环）。
